\begin{thebibliography}{99}

\bibitem[Afriat(1967)]{afriat1967}
Afriat, S. N. (1967).
The construction of utility functions from expenditure data.
\emph{International Economic Review} 8, 67--77.

\bibitem[Bellman(1957)]{bellman1957}
Bellman, R. (1957).
\emph{Dynamic Programming}.
Princeton, NJ: Princeton University Press.

\bibitem[Elaydi(2005)]{elaydi2005}
Elaydi, S. (2005).
\emph{An Introduction to Difference Equations}, 3rd ed.
New York: Springer.

\bibitem[Bhattacharya(2015)]{bhattacharya2015}
Bhattacharya, D. (2015).
Nonparametric welfare analysis for discrete choice.
\emph{Econometrica} 83, 617--649.

\bibitem[Brown and Calsamiglia(2007)]{browncalsamiglia2007}
Brown, D. J. and C. Calsamiglia (2007).
The nonparametric approach to applied welfare analysis.
\emph{Economic Theory} 31, 183--188.

\bibitem[Castillo and Freer(2023)]{castillofreer2023}
Castillo, M. and M. Freer (2023).
A general revealed preference test for quasilinear preferences: theory and experiments.
\emph{Experimental Economics} 26, 673--696.

\bibitem[Chambers, Echenique, and Lambert(2021)]{chambersetal2021}
Chambers, C. P., F. Echenique, and N. S. Lambert (2021).
Recovering preferences from finite data.
\emph{Econometrica} 89, 1633--1664.

\bibitem[Debreu(1954)]{debreu1954}
Debreu, G. (1954).
Representation of a preference ordering by a numerical function.
In R. M. Thrall, C. H. Coombs, and R. L. Davis (eds.), \emph{Decision Processes}, 159--165. New York: Wiley.

\bibitem[Doignon and Falmagne(1974)]{doignonfalmagne1974}
Doignon, J. P. and J. C. Falmagne (1974).
Difference measurement and simple scalability with restricted solvability.
\emph{Journal of Mathematical Psychology} 11, 473--499.

\bibitem[Fishburn(2001)]{fishburn2001}
Fishburn, P. C. (2001).
Cancellation conditions for finite two-dimensional additive measurement.
\emph{Journal of Mathematical Psychology} 45, 2--26.

\bibitem[Fishburn, Marcus-Roberts, and Roberts(1988)]{fishburnetal1988}
Fishburn, P. C., H. M. Marcus-Roberts, and F. S. Roberts (1988).
Unique finite difference measurement.
\emph{SIAM Journal on Discrete Mathematics} 1, 334--354.

\bibitem[Fuchs-Seliger(1989)]{fuchsseliger1989}
Fuchs-Seliger, S. (1989).
Money-metric utility functions in the theory of revealed preference.
\emph{Mathematical Social Sciences} 18, 199--210.

\bibitem[Fuchs-Seliger(1990)]{fuchsseliger1990}
Fuchs-Seliger, S. (1990).
An axiomatic approach to compensated demand.
\emph{Journal of Economic Theory} 52, 111--122.

\bibitem[Gorno(2019)]{gorno2019}
Gorno, L. (2019).
Revealed preference and identification.
\emph{Journal of Economic Theory} 183, 698--739.

\bibitem[Kraft, Pratt, and Seidenberg(1959)]{kraftprattseidenberg1959}
Kraft, C. H., J. W. Pratt, and A. Seidenberg (1959).
Intuitive probability on finite sets.
\emph{Annals of Mathematical Statistics} 30, 408--419.

\bibitem[Kuznetsov(2004)]{kuznetsov2004}
Kuznetsov, Y. A. (2004).
\emph{Elements of Applied Bifurcation Theory}, 3rd ed.
New York: Springer.

\bibitem[Makowski and Ostroy(2013)]{makowskiostroy2013}
Makowski, L. and J. M. Ostroy (2013).
From revealed preference to preference revelation.
\emph{Journal of Mathematical Economics} 49, 71--81.

\bibitem[Nocke and Schutz(2017)]{nockeschutz2017}
Nocke, V. and N. Schutz (2017).
Quasi-linear integrability.
\emph{Journal of Economic Theory} 169, 603--628.

\bibitem[Rochet(1987)]{rochet1987}
Rochet, J.-C. (1987).
A necessary and sufficient condition for rationalizability in a quasi-linear context.
\emph{Journal of Mathematical Economics} 16, 191--200.

\bibitem[Samuelson(1938)]{samuelson1938}
Samuelson, P. A. (1938).
A note on the pure theory of consumer's behaviour.
\emph{Economica} 5, 61--71.

\bibitem[Samuelson(1948)]{samuelson1948}
Samuelson, P. A. (1948).
Consumption theory in terms of revealed preference.
\emph{Economica} 15, 243--253.

\bibitem[Samuelson(1950)]{samuelson1950}
Samuelson, P. A. (1950).
The problem of integrability in utility theory.
\emph{Economica} 17, 355--385.

\bibitem[Scott(1964)]{scott1964}
Scott, D. (1964).
Measurement structures and linear inequalities.
\emph{Journal of Mathematical Psychology} 1, 233--247.

\bibitem[Stokey, Lucas, and Prescott(1989)]{stokeylucasprescott1989}
Stokey, N. L., R. E. Lucas, Jr., and E. C. Prescott (1989).
\emph{Recursive Methods in Economic Dynamics}.
Cambridge, MA: Harvard University Press.

\end{thebibliography}

\end{document}
